%% 
%% Copyright 2007, 2008, 2009 Elsevier Ltd
%% 
%% This file is part of the 'Elsarticle Bundle'.
%% ---------------------------------------------
%% 
%% It may be distributed under the conditions of the LaTeX Project Public
%% License, either version 1.2 of this license or (at your option) any
%% later version.  The latest version of this license is in
%%    http://www.latex-project.org/lppl.txt
%% and version 1.2 or later is part of all distributions of LaTeX
%% version 1999/12/01 or later.
%% 
%% The list of all files belonging to the 'Elsarticle Bundle' is
%% given in the file `manifest.txt'.
%% 

%% Template article for Elsevier's document class `elsarticle'
%% with numbered style bibliographic references
%% SP 2008/03/01

\documentclass[preprint,12pt, a4paper]{elsarticle}

%% Use the option review to obtain double line spacing
%% \documentclass[authoryear,preprint,review,12pt]{elsarticle}

%% For including figures, graphicx.sty has been loaded in
%% elsarticle.cls. If you prefer to use the old commands
%% please give \usepackage{epsfig}

%% The amssymb package provides various useful mathematical symbols
\usepackage{amssymb}
\usepackage{hyperref}
\setlength{\parindent}{0pt}
%% The amsthm package provides extended theorem environments
%% \usepackage{amsthm}

%% The lineno packages adds line numbers. Start line numbering with
%% \begin{linenumbers}, end it with \end{linenumbers}. Or switch it on
%% for the whole article with \linenumbers.
%\usepackage{lineno}

\journal{SoftwareX}

\begin{document}
\renewcommand{\labelenumii}{\arabic{enumi}.\arabic{enumii}}

\begin{frontmatter}
%--- INSTRUCTIONS TO BE DELETED OR COMMENTED BEFORE SUBMISSION 

  {\large\textbf{SoftwareX article template Version 4 (November 2023)}}
  Before you complete this template, a few important points to note:
  \begin{itemize}
\item	This template is for an original SoftwareX article. If you are submitting an update to a software article that has already been published, please use the Software Update Template found on the Guide for Authors page.
\item	The format of a software article is very different to a traditional research article. To help you write yours, we have created this template. We will only consider software articles submitted using this template.
\item	It is mandatory to publicly share the code/software referred to in your software article. You’ll find information on our software sharing criteria in the SoftwareX \href{https://www.elsevier.com/journals/softwarex/2352-7110/guide-for-authors}{Guide for Authors}.  
\item	It’s important to consult the \href{https://www.elsevier.com/journals/softwarex/2352-7110/guide-for-authors}{Guide for Authors} when preparing your manuscript; it highlights mandatory requirements and is packed with useful advice.
\end{itemize}
%
Still got questions?
Email our editorial team at softwarex@elsevier.com.

\textbf{Now you are ready to fill in the template below. As you complete each section, please carefully read the associated instructions. All sections are mandatory, unless marked optional.
Once you have completed the template, delete these instructions. In addition, please delete the instructions in the template (the text written in italics).}
%--- END OF INSTRUCTIONS TO BE DELETED OR COMMENTED BEFORE SUBMISSION 
 
%% Title, authors and addresses

%% use the tnoteref command within \title for footnotes;
%% use the tnotetext command for theassociated footnote;
%% use the fnref command within \author or \address for footnotes;
%% use the fntext command for theassociated footnote;
%% use the corref command within \author for corresponding author footnotes;
%% use the cortext command for theassociated footnote;
%% use the ead command for the email address,
%% and the form \ead[url] for the home page:
%% \title{Title\tnoteref{label1}}
%% \tnotetext[label1]{}
%% \author{Name\corref{cor1}\fnref{label2}}
%% \ead{email address}
%% \ead[url]{home page}
%% \fntext[label2]{}
%% \cortext[cor1]{}
%% \address{Address\fnref{label3}}
%% \fntext[label3]{}


%% use optional labels to link authors explicitly to addresses:
%% \author[label1,label2]{}
%% \affiliation[label1]{organization={},
%%             addressline={},
%%             city={},
%%             postcode={},
%%             state={},
%%             country={}}
%%
%% \affiliation[label2]{organization={},
%%             addressline={},
%%             city={},
%%             postcode={},
%%             state={},
%%             country={}}

            




\title{EnerHabitat: A Python package and GUI to solve 1D  time-dependent heat transfer on building components} %% Article title

%% use optional labels to link authors explicitly to addresses:
%% \author[label1,label2]{}
%% \address[label1]{}
%% \address[label2]{}


\author{Fernando Rodríguez Calderón} %% Author name

%% Author affiliation
\affiliation{organization={IPN},%Department and Organization
            addressline={Direccion poli}, 
            city={Ciudad de Mexico},
            postcode={12345}, 
            state={CdMx},
            country={Mexico}}



\author{Guillermo Barrios} %% Author name
\author{Miriam Cruz Salas} %% Author name

\author{Guadalupe Huelsz} %% Author name
\author{Jorge Rojas} %% Author name

%% Author affiliation
\affiliation{organization={Instituto de Energías Renovables, UNAM},%Department and Organization
            addressline={}, 
            city={Temixco},
            postcode={62580}, 
            state={Morelos},
            country={Mexico}}

% \author[label1]{A. Author1}
% \author[label2]{A. Author2}
% \address[label1]{Author1's affiliation, address, email}
% \address[label2]{Author2's affiliation, address, email}

\begin{abstract}
%% Text of abstract 
\textit{Ca. 100 words. The abstract is followed by a maximum of six keywords
and some mandatory and optional metadata.}

\textit{Your main body of text (sections 1-5 below) should be a maximum 6 pages in total (excluding metadata, tables, figures, references) with a 3000-word limit (we ask that more priority is placed on the word limit versus the page count). Though we strictly insist on the author following the template, in exceptional circumstances, we can be flexible with the page numbers and word limit. In such cases, it should be discussed with the managing editor or publisher prior to submission. All queries regarding the same can be reached at softwarex@elsevier.com.}
\end{abstract}

\begin{keyword}
%% keywords here, in the form: keyword \sep keyword
time-dependent heat transfer \sep building envelope \sep thermal performance \sep Shiny for Python

%% PACS codes here, in the form: \PACS code \sep code

%% MSC codes here, in the form: \MSC code \sep code
%% or \MSC[2008] code \sep code (2000 is the default)

\end{keyword}

\end{frontmatter}

%\linenumbers

\section*{Metadata}
\label{}
\textit{The ancillary data table~\ref{codeMetadata} is required for the sub-version of the codebase. Please replace the italicized text in the right column with the correct information about your current code and leave the left column untouched.}

\begin{table}[!h]
\begin{tabular}{|l|p{6.5cm}|p{6.5cm}|}
\hline
\textbf{Nr.} & \textbf{Code metadata description} & \textbf{Metadata} \\
\hline
C1 & Current code version & For example v42 \\
\hline
C2 & Permanent link to code/repository used for this code version & For example: \url{https://github.com/mozart/mozart2} \\
\hline
C3  & Permanent link to Reproducible Capsule & For example: \url{https://codeocean.com/capsule/0270963/tree/v1}\\
\hline
C4 & Legal Code License   & All software and code must be released under one of the pre-approved licenses listed in the \href{https://www.elsevier.com/journals/softwarex/2352-7110/guide-for-authors}{Guide for Authors}, such as Apache License, GNU General Public License (GPL) or MIT License. Write the name of the license you’ve chosen here. \\
\hline
C5 & Code versioning system used & For example: svn, git, mercurial,
                                   etc. (put none if none is used) \\
\hline
C6 & Software code languages, tools, and services used & For example: C++, python, r, MPI, OpenCL, etc. \\
\hline
C7 & Compilation requirements, operating environments \& dependencies & \\
\hline
C8 & If available Link to developer documentation/manual & For example: \url{http://mozart.github.io/documentation/} \\
\hline
C9 & Support email for questions & \\
\hline
\end{tabular}
\caption{Code metadata (mandatory)}
\label{codeMetadata} 
\end{table}

\textit{Optionally, you can provide information about the current executable
software version filling in the left column of
Table~\ref{executabelMetadata}. Please leave the first column as it is.}

\begin{table}[!h]
\begin{tabular}{|l|p{6.5cm}|p{6.5cm}|}
\hline
\textbf{Nr.} & \textbf{(Executable) software metadata description} & \textbf{Please fill in this column} \\
\hline
S1 & Current software version & For example 1.1, 2.4 etc. \\
\hline
S2 & Permanent link to executables of this version  & For example: \url{https://github.com/combogenomics/DuctApe/releases/tag/DuctApe-0.16.4} \\
\hline
S3  & Permanent link to Reproducible Capsule & \\
\hline
S4 & Legal Software License & List one of the approved licenses \\
\hline
S5 & Computing platforms/Operating Systems & For example Android, BSD, iOS, Linux, OS X, Microsoft Windows, Unix-like , IBM z/OS, distributed/web based etc. \\
\hline
S6 & Installation requirements \& dependencies & \\
\hline
S7 & If available, link to user manual - if formally published include a reference to the publication in the reference list & For example: \url{http://mozart.github.io/documentation/} \\
\hline
S8 & Support email for questions & \\
\hline
\end{tabular}
\caption{Software metadata (optional)}
\label{executabelMetadata} 
\end{table}


\section{Motivation and significance}


The thermal behavior of opaque building envelopes plays a central role in the energy performance and indoor comfort of buildings, particularly in warm climates and in free-running dwellings where mechanical cooling is absent or unaffordable. In many regulatory frameworks and design tools, envelope components are still characterized using steady-state indicators such as U-values or R-values. These indicators neglect thermal capacity and the dynamic interaction between solar gains, outdoor temperature cycles and internal heat storage, which can lead to misleading design decisions and to regulations that are not well suited for lightweight or non-homogeneous constructions. This limitation is especially critical in social housing and in regions with strong diurnal temperature swings and high solar irradiation.

EnerHabitat was originally developed to address this gap through a time-dependent numerical model for single envelope components (walls and roofs), capable of handling both homogeneous multilayer systems and systems that include a non-homogeneous layer with air cavities, under either air-conditioned or non air-conditioned indoor conditions for several cities of Mexico.\citep{barrios2016enerhabitat} The models were validated against EnergyPlus for homogeneous assemblies and against hot-box experiments for hollow concrete block walls, demonstrating good agreement in indoor air temperatures, decrement factors, lag times and heat fluxes. However, the initial implementation was delivered as a closed, web-based tool with  a pre-defined climate database limited possibilities for automation, integration into scientific workflows or community-driven extension and available only through an online platform and without access to the source.

The open-source version of Ener-Habitat aims to make these validated models fully transparent, reproducible and extensible. By releasing the core algorithms and data structures under an open license, the software directly addresses several scientific problems:

\begin{itemize}
\item It provides an accessible reference implementation of the time-dependent one-dimensional  heat transfer in multilayer envelope systems with convective boundary conditions.
\item It allows designers or researchers to systematically explore how envelope material choices, layer ordering and non-homogeneous geometries affect dynamic metrics such as decrement factor, time lag, indoor temperature swings, heating/cooling loads and discomfort degree-hours in both air-conditioned and free-running scenarios.
\item It enables the development and testing of new evaluation criteria and regulatory indicators for warm-climate envelopes, beyond steady-state U-values, and therefore supports the design of more appropriate building energy standards.
\end{itemize}

The original Ener-Habitat models have already been used to highlight differences in wall/roof performance between air-conditioned and non air-conditioned rooms, and to propose thermal performance parameters tailored to free-running buildings.\citep{barrios2011wall,barrios2012envelope} By exposing these models as open-source software, we lower the barrier for other researchers to reuse and extend them in comparative studies, parametric analyses and optimisation workflows (for example, coupling envelope simulations with life-cycle assessment, cost optimisation or climate-change scenarios). In teaching, the open-source implementation can be integrated into Jupyter notebooks and classroom exercises, enabling students to directly manipulate envelope configurations, climates and indoor conditions, and to visualise the resulting thermal response.

From the user perspective, the experimental setting is straightforward. A typical workflow consists of: (i) selecting or importing climate data for a given location and period; (ii) defining one or more constructive systems as a sequence of layers, each with thickness and thermo-physical properties, possibly including one non-homogeneous block layer with specified geometry; (iii) choosing the envelope position (wall or roof), inclination and orientation; (iv) specifying whether the room is air-conditioned (with a prescribed neutral temperature) or free-running; and (v) running the simulation to obtain time series of outdoor air, sol–air, surface and indoor air temperatures, as well as summary indicators such as transmitted energy, decrement factor, lag time and the dataset related to the evaluation. In the open-source implementation, this workflow can be executed either through a graphical interface for interactive exploration or via python workflow  suitable for large parametric studies.

Several families of tools are related to EnerHabitat. Whole-building simulation programs such as EnergyPlus, TRNSYS, ESP-r, BSim, DeST, IESVe and others provide detailed dynamic modelling of entire buildings, but at the cost of steep learning curves and complex input structures, and they are not optimised for rapid exploration of single envelope components or for systematically varying non-homogeneous block geometries. Steady-state or quasi-steady tools and calculators (e.g. PHPP, Lesokai, Construction R-value Calculator) focus on U-values and simple indicators, and thus cannot capture the dynamic attenuation and time lag of heat waves through walls and roofs. Two- and three-dimensional conduction tools such as THERM, AnTherm, Heat2 and Heat3 can represent thermal bridges and non-homogeneous details, but they are typically used in steady-state or harmonic regimes and do not provide the envelope-oriented performance metrics needed by designers in warm climates. 
% \citep[e.g.][]{kossecka2002influence,vijayalakshmi2006thermal,ozel2007optimum,gao2004reduced,alhazmy2006analysis,alkhameis2011coupled,alsanea2012effect,alsanea2013effect,zhang2014impact} 

Within this landscape, the open-source EnerHabitat occupies a distinctive niche: it packages validated one-dimensional transient model in a lightweight, envelope-focused tool; it exposes design-relevant outputs for both air-conditioned and free-running buildings; and, by being open-source, it invites scrutiny, extension and integration by the building science community. 



\section{Software description}




EnerHabitat is currently distributed as an open-source ecosystem composed of three tightly coupled elements: (i) a core Python package (\texttt{enerhabitat}) published on the Python Package Index and hosted on GitHub; (ii) a graphical user interface implemented with Shiny for Python in a separate repository; and (iii) a companion validation repository that contains input data and notebooks to reproduce the numerical and experimental validation cases documented in previous work. Together, these components provide a reproducible, scriptable and user-friendly environment for the dynamic evaluation of opaque envelope systems based on standard EPW weather files.

\subsection{Software architecture}

At a high level, the architecture is organised in layers that follow the physical workflow from climate data to indoor temperature response:

\begin{enumerate}
  \item \textbf{Core numerical engine (\texttt{enerhabitat} package).} The Python package encapsulates the data structures and numerical routines needed to (a) read and pre-process EPW weather files, (b) compute a representative ``mean day'' with sub-hourly resolution, (c) derive sol--air boundary conditions, and (d) solve the transient one-dimensional heat conduction problem through a multilayer constructive system. The main user-facing functions are:
  \begin{itemize}
    \item \texttt{meanDay}, which reads an EPW file and returns a \texttt{pandas} data frame containing, for each time step, solar position variables and meteorological quantities such as ambient temperature and the global, direct and diffuse components of irradiance.
    \item \texttt{Tsa}, which takes the output of \texttt{meanDay} and computes the sol--air temperature and surface irradiance for a given surface tilt, azimuth and solar absorptance.
    \item \texttt{solveCS}, which takes a constructive system defined as an ordered list of (material, thickness) pairs and solves the transient conduction problem for the chosen sol--air boundary condition, returning the indoor air temperature time series for the fictitious room associated with that envelope.
  \end{itemize}
  Material properties are stored in an external \texttt{INI} configuration file (\texttt{materials.ini}), where each material is defined by its thermal conductivity (\(k\)), density (\(\rho\)) and specific heat capacity (\(c\)). A small set of global parameters controls the numerical resolution and boundary conditions, including the characteristic room length \texttt{La}, the number of spatial elements \texttt{Nx}, the internal and external convective heat transfer coefficients (\texttt{hi} and \texttt{ho}), and the time step \texttt{dt}. Internally, the core uses \texttt{numpy} arrays for vectorised operations and \texttt{numba} to JIT-compile the time-stepping loops, while \texttt{pvlib} is used for solar geometry and irradiance calculations.
  
  \item \textbf{Shiny for Python web application.} The graphical interface (\texttt{enerhabitat\_webapp} repository) is a thin presentation layer built with Shiny for Python on top of the core package. It organises the workflow into interactive panels where the user can (a) upload or select EPW files, (b) choose or edit the materials and construct the wall or roof assembly, (c) specify geometric parameters (tilt, azimuth) and indoor boundary conditions (air-conditioned vs.\ free-running room), and (d) execute the simulation and visualise the resulting time series and summary indicators. All computational tasks are delegated to the \texttt{enerhabitat} package, so that the web application remains stateless and easy to maintain: interface callbacks only collect user input, call \texttt{meanDay}, \texttt{Tsa} and \texttt{solveCS}, and then map the returned data frames to plots and tables.
  
  \item \textbf{Validation and reproducibility repository.} The \texttt{validacion\_enerhabitat} repository provides a curated set of input data and scripts for reproducing the validation procedures described in the methodological work behind EnerHabitat. It includes EPW files, EnergyPlus/OpenStudio models (\texttt{.idf}, \texttt{.osm}), measurement data from hot-box experiments, and Jupyter notebooks that call the \texttt{enerhabitat} package to compare its results against full building simulations and experimental time series. This repository functions both as an executable appendix to the scientific publications and as a set of advanced examples for users who want to integrate EnerHabitat into larger simulation workflows.
\end{enumerate}

Figure~\ref{fig:architecture} (to be added by the authors) can schematically summarise these components and the data flow: EPW files are processed by \texttt{meanDay}, the resulting time series feed \texttt{Tsa}, and the sol--air boundary condition is passed to \texttt{solveCS}, with outputs consumed either programmatically (in Python scripts or notebooks) or through the Shiny web interface.

\subsection{Software functionalities}

From the end-user perspective, the main functionalities can be grouped into four categories:

\begin{enumerate}
  \item \textbf{Meteorological pre-processing.} EnerHabitat reads EPW weather files and condenses the selected period (typically a year) into a synthetic mean day with a user-defined temporal resolution. The output data frame includes solar position (zenith, elevation, azimuth, equation of time), ambient temperature and the three irradiance components (global horizontal, direct normal and diffuse horizontal), which can be directly inspected or reused in other analyses.

  \item \textbf{Boundary-condition construction (sol--air temperature).} Given the mean-day meteorology, the \texttt{Tsa} function computes the sol--air temperature and the surface irradiance for a specified surface tilt, azimuth and solar absorptance. This step encapsulates the radiative and convective exchange at the external surface, providing a compact boundary condition that drives the subsequent conduction calculation. Users can rapidly explore the effect of orientation, colour (through absorptance) and slope (wall vs.\ roof) on the external thermal forcing.

  \item \textbf{Dynamic simulation of multilayer constructive systems.} The \texttt{solveCS} function implements the transient, one-dimensional heat-conduction model for a wall or roof assembly defined as a sequence of layers. Each layer refers to a material in the configuration file, which ensures that material libraries can be maintained independently of the code. The function returns at least the indoor air temperature time series associated with the fictitious room behind the envelope; from this output, users can derive dynamic performance indicators such as decrement factor, time lag, daily temperature amplitude, cumulative heat flux or temperature-based comfort metrics. The numerical resolution and convective coefficients can be tuned via global parameters, allowing the trade-off between accuracy and computational cost to be explicitly controlled.

  \item \textbf{Interactive exploration and visualisation.} The Shiny for Python interface exposes the above functionalities to non-programming users. Typical operations include selecting a climate file from a list, choosing a wall or roof type from a catalogue of constructive systems, sketching a custom assembly by adding or removing layers, and immediately seeing the impact on indoor temperature curves and summary indicators. Results can be exported (e.g.\ as CSV files) for inclusion in reports or further analysis. For expert users, the validation repository offers more advanced scripts and notebooks that demonstrate how to embed EnerHabitat into larger simulation chains, for example when comparing its predictions with EnergyPlus or with experimental data.
\end{enumerate}

\subsection{Sample code snippets analysis}

Listing~\ref{lst:basic-usage} illustrates a minimal Python script that uses the core \texttt{enerhabitat} package to evaluate a wall assembly for a given climate file. The same sequence of operations---EPW pre-processing, sol--air computation and constructive-system solution---underlies the graphical web interface.

\begin{verbatim}
import enerhabitat as eh

# 1. Read the EPW file and build the mean-day meteorology
mean_day = eh.meanDay(epw_file="epw/Temixco_MEX.epw")

# 2. Build the sol-air boundary condition for an east-facing wall
tsa = eh.Tsa(
    mean_day,
    solar_absortance=0.8,  # medium-dark surface
    surface_tilt=90.0,     # vertical surface
    surface_azimuth=90.0   # 0 = North, 90 = East
)

# 3. Define the constructive system from exterior to interior
system = [
    ("exterior_plaster", 0.015),
    ("hollow_concrete_block", 0.12),
    ("interior_plaster", 0.015),
]

# Optional: adjust numerical and boundary-condition parameters
eh.La = 2.5    # characteristic room length [m]
eh.Nx = 120    # number of spatial elements
eh.ho = 13.0   # outside convection coefficient [W/m²K]
eh.hi = 8.6    # inside convection coefficient [W/m²K]
eh.dt = 60.0   # time step [s]

# 4. Solve the transient 1D heat conduction problem
interior = eh.solveCS(system, tsa)

# 'interior' is a pandas Series or DataFrame with indoor temperature vs. time
\end{verbatim}

In this example, the user only needs to provide three pieces of information: an EPW file (\texttt{Temixco\_MEX.epw}), the optical and geometric properties of the envelope surface, and a list describing the wall layers. All other parameters have sensible defaults but can be overridden if needed. The returned indoor temperature can be directly plotted, compared between alternative assemblies or post-processed to obtain comfort or performance indicators. Within the Shiny interface, the same logic is followed: each widget corresponds to one argument in the underlying functions, and each button click triggers a call to \texttt{meanDay}, \texttt{Tsa} and \texttt{solveCS} with the values chosen by the user.




To illustrate the main capabilities of EnerHabitat, we describe two representative use cases.

\subsection{Example 1: Comparing wall assemblies for a free-running social dwelling}

In the first example, EnerHabitat is used to compare two alternative wall assemblies for a naturally ventilated social housing unit located in a warm sub-humid climate (e.g.\ Temixco, Morelos). The workflow is as follows:

\begin{enumerate}
  \item The user selects the local EPW file and specifies a representative surface orientation (for instance, an east-facing bedroom wall).
  \item Two constructive systems are defined: a baseline wall (e.g.\ hollow concrete blocks with thin interior and exterior plaster layers) and an improved wall including an insulating layer or higher thermal capacity materials.
  \item For each assembly, \texttt{meanDay} and \texttt{Tsa} are called to compute the mean-day boundary conditions, and \texttt{solveCS} is used to obtain the indoor temperature time series of the fictitious room behind the wall.
  \item The resulting indoor temperature profiles are plotted for the 24~h period and compared in terms of maximum and minimum temperatures, daily amplitude, time lag between outdoor and indoor peaks, and the number of hours above a chosen thermal comfort threshold.
\end{enumerate}

Using the Python API, these steps can be implemented in a short script or notebook with a loop over the two assemblies, while the Shiny interface allows the same comparison to be performed interactively by selecting each constructive system in turn. In both cases, EnerHabitat enables rapid, dynamic comparisons that would be more cumbersome to carry out in a full whole-building simulation environment.

\subsection{Example 2: Reproducing published validation cases}

In the second example, a researcher uses the \texttt{validacion\_enerhabitat} repository to reproduce one of the published validation cases for EnerHabitat. The repository provides the necessary EPW files, EnergyPlus/OpenStudio models and measurement data, along with Jupyter notebooks that:

\begin{enumerate}
  \item Construct the corresponding wall or roof assembly in \texttt{enerhabitat} using the same material properties and layer thicknesses as in the reference model or experiment.
  \item Run the simulation with \texttt{meanDay}, \texttt{Tsa} and \texttt{solveCS}.
  \item Post-process the results to compute metrics such as indoor air temperature, decrement factor and lag time, as well as heat fluxes at the interior surface.
  \item Overlay the EnerHabitat predictions with the EnergyPlus results or experimental measurements, reproducing the figures and error metrics reported in the original validation work.
\end{enumerate}

This example highlights how EnerHabitat, together with its validation repository, supports open and reproducible research: all input files, intermediate data and analysis scripts are version-controlled and publicly available, and the numerical core is the same in the command-line, notebook and web-interface workflows.


\section{Impact}

EnerHabitat targets a long-standing gap between detailed whole-building simulators and simple steady-state calculators: the need for a lightweight, envelope-centric tool that captures the dynamic behaviour of walls and roofs under realistic outdoor forcing, especially in warm climates and free-running dwellings.\citep{barrios2016enerhabitat} By re-implementing the original online tool as an open-source Python package with a Shiny for Python interface, EnerHabitat becomes easier to inspect, extend and couple to other scientific workflows, which substantially increases its potential impact within building science and energy policy.

\paragraph{New research questions enabled.}

The open-source implementation makes it feasible to pose and answer research questions that were difficult to address with the original closed web application or with general-purpose whole-building tools. Examples include:

\begin{itemize}
\item Systematic exploration of how low-cost material substitutions, layer ordering and non-homogeneous block geometries affect dynamic metrics such as decrement factor, time lag and indoor temperature swings in free-running rooms.
\item Development and testing of alternative envelope performance indicators beyond steady-state U-value, tailored to hot or warm-humid climates and naturally ventilated housing.\citep{barrios2011wall,barrios2012envelope}
\item Coupling envelope-level simulations with life cycle assessment, cost optimisation or climate-change scenarios, by embedding EnerHabitat into Python-based pipelines for parametric studies.
\item Comparative studies between reduced-order envelope models and detailed whole-building simulations (e.g.\ EnergyPlus, TRNSYS), using the validation repository as a starting point.
\end{itemize}

Because the numerical core is now a normal Python package, these questions can be explored programmatically (for example, by sampling large design spaces or climate scenarios) rather than one configuration at a time through a web form.

\paragraph{Improving existing research workflows.}

Even for research questions that were already being pursued with the original Ener-Habitat tool and with other simulators, the new implementation improves the process in several ways. First, it provides a transparent and reproducible reference implementation of the validated one- and two-dimensional heat transfer models documented in the original paper, making it straightforward for other groups to replicate and build upon previous results.\citep{barrios2016enerhabitat} Second, it lowers the barrier to integrating envelope simulations into existing Python ecosystems, where researchers already use \texttt{numpy}, \texttt{pandas} and \texttt{pvlib} for data handling and solar geometry.

Typical examples include: (i) replacing ad-hoc spreadsheets or bespoke code with a tested, version-controlled package; (ii) embedding EnerHabitat calls inside Jupyter notebooks used to generate figures and tables for publications; and (iii) using the validation repository as executable documentation to verify that local installations reproduce the published comparison with EnergyPlus and hot-box experiments. In this sense, the software acts as a “bridge” between detailed building physics and modern, open scientific computing practices.

\paragraph{Changes in daily practice.}

For its original user community—building energy researchers, teachers and students at the Instituto de Energías Renovables (IER–UNAM) and collaborating groups—EnerHabitat has already changed daily practice in several ways. It is used as a teaching tool in courses on energy in buildings and numerical methods, where students learn to go from EPW weather data to time-varying indoor temperatures and dynamic performance indicators using a single, coherent set of functions, and more recently through the graphical Shiny interface.\footnote{See, for example, the author’s academic profile and workshop descriptions, where EnerHabitat is listed as a free-access tool for evaluating wall and roof performance and as part of sustainability-oriented educational projects.([ier.unam.mx][1])}

For design-oriented studies, EnerHabitat is used as a rapid screening tool for envelope options before investing effort in full whole-building models. Practitioners can test “what if” scenarios—such as changing the colour (solar absorptance), thickness or internal configuration of a wall—and immediately see the impact on indoor temperature amplitude and discomfort hours. This has led to a more systematic use of dynamic envelope criteria in early-stage design and in applied research projects focused on Mexican social housing.

\paragraph{Adoption and dissemination.}

The original web-based Ener-Habitat has been available as a free online service for nearly a decade and has been cited and used in graduate work and research on Mexican housing and regulations, where it appears explicitly as software supporting the analysis of envelope performance.\citep{barrios2016enerhabitat} It is also referenced as an accessible tool for evaluating walls and roofs in institutional documents and theses at UNAM.([Tesi UNAM Documentos][2])

The new open-source implementation is more recent: the core package is now hosted on GitHub and released on the Python Package Index (PyPI), with binary archives mirrored on public university mirrors, which ensures straightforward installation for Python users worldwide.([GitHub][3]) A dedicated repository hosts the Shiny for Python web application, and a separate validation repository packages the input data and notebooks that reproduce the published validation cases.([GitHub][4])

At the time of writing, most users are still concentrated in the intended core audience—researchers, graduate students and practitioners working on building envelopes in warm-climate regions—but the packaging and documentation are designed to facilitate wider adoption. The fact that EnerHabitat is now installable via \texttt{pip} and can be integrated into existing Python projects significantly lowers the overhead for external users compared to the original web-only deployment.

\paragraph{Commercial use and spin-offs.}

EnerHabitat is currently used primarily in academic, educational and non-profit contexts. We are not aware of direct commercial deployments or spin-off companies built specifically around the software. This is partly by design: the tool is released under a permissive open-source licence (MIT) to encourage its uptake as an infrastructural component within existing design and consulting workflows rather than as a closed commercial product.([GitHub][3])

Nevertheless, the open-source distribution creates clear opportunities for commercial actors. Engineering firms and independent consultants can embed EnerHabitat into their internal toolchains without licensing barriers, for example to pre-screen envelope options, document envelope performance in reports, or support compliance studies related to dynamic envelope criteria. Future collaborations with industry partners may lead to domain-specific interfaces or packaged services (e.g.\ web dashboards for social housing agencies) that build directly on the EnerHabitat core.

\section{Conclusions}

EnerHabitat provides an open, extensible implementation of a family of validated numerical models for the dynamic thermal evaluation of opaque envelope components, together with a user-friendly Shiny for Python interface and an executable validation repository. By focusing on single-component, wall and roof assemblies and by embracing an open-source Python ecosystem, the software occupies a niche between detailed whole-building simulators and simplified steady-state tools, enabling rapid yet physically grounded exploration of envelope design options, especially in warm climates and free-running dwellings.

From a scientific perspective, EnerHabitat consolidates previous methodological work on one- and two-dimensional heat transfer models for homogeneous and non-homogeneous constructive systems into reusable code, thus supporting reproducibility and further methodological extensions.\citep{barrios2016enerhabitat,barrios2011wall,barrios2012envelope} From a practical standpoint, it lowers the barrier for researchers, students and practitioners to incorporate dynamic envelope criteria into their daily workflows, whether through interactive use of the Shiny interface or through scripted analyses in Jupyter notebooks and Python pipelines.

Future developments will focus on several directions. On the modelling side, we plan to extend the library of non-homogeneous configurations beyond the current block geometries, and to include additional output metrics related to comfort, overheating risk and life-cycle performance. On the software-engineering side, we aim to strengthen the testing and documentation infrastructure, provide containerised deployments, and explore tighter coupling with whole-building simulators and optimisation frameworks. Finally, by maintaining the project as a community-driven, open-source effort, we hope that EnerHabitat will continue to evolve as a shared platform for research, teaching and practice in building envelope design for warm-climate, low-energy housing.




\section*{Acknowledgements}
\label{}
\textit{Optional. You can use this section to acknowledge colleagues who don’t qualify as a co-author but helped you in some way. }


%% The Appendices part is started with the command \appendix;
%% appendix sections are then done as normal sections
%% \appendix

%% \section{}
%% \label{}

%% References:
%% If you have bibdatabase file and want bibtex to generate the
%% bibitems, please use
%%
%%  \bibliographystyle{elsarticle-num} 
%%  \bibliography{<your bibdatabase>}

%% else use the following coding to input the bibitems directly in the
%% TeX file.

\begin{thebibliography}{00}

%% \bibitem{label}
%% Text of bibliographic item

\bibitem{} Use this style of ordering. References in-text should also use a similar style.

\end{thebibliography}

\textit{If the software repository you used supplied a DOI or another
Persistent IDentifier (PID), please add a reference for your software
here. For more guidance on software citation, please see our guide for
authors or \href{https://f1000research.com/articles/9-1257/v2}{this
  article on the essentials of software citation by FORCE 11}, of
which Elsevier is a member.}

\large{\textbf{Reminder: Before you submit, please delete all 
the instructions in this document, 
including this paragraph. 
Thank you!}}




\end{document}
\endinput
%%
%% End of file `SoftwareX_article_template.tex'.

%%% Local Variables:
%%% mode: latex
%%% TeX-master: t
%%% End:
